% 18 November 2026. Drafted for Edward Baker with OpenAI Codex assistance.
% Actual model and effort unavailable; authoring-helper settings inherited.
\section{Retained forcing and the exact backward return allowance}
\label{sec:retained-forcing}

A whole-time forcing estimate carries more information than an initial
interaction or a constant upper bound for the defect. We now keep that
information in two scalar comparisons and then identify exactly what is
lost when a forward evolution is compressed back to its initial-$M$
backward horizon. The derivations are from \cite{BakerCM32,BakerCM34},
with the internal assessments \cite{BakerCMReview29,BakerCMReview30}.
The numerical anchor tests are deferred to Appendix~\ref{app:compact-example}.
All statements in this section are analytical and independent of those
supporting values.

Work first in physical backward time $u=t_*-t$ on a counted pair tube
where the complete local equations hold. Write
\begin{equation}
 D_u=-2+4DM,\qquad M_u=-2M^2+\mathcal F,\qquad
 \mathcal F=12S_2+16DS_3,\qquad K=2D^2\mathcal F.
 \label{rf:pair-system}
\end{equation}
The pair is real and distinct while $D>0$. Its exterior may be nonreal:
the complete signed quantities in \eqref{rf:pair-system} are then
acquired through Section~\ref{sec:compact}, rather than by omitting an
uncounted tail. With $V=1-2DM$, the identities
\begin{equation}
 D_u=-2V,\qquad V_u=\frac{1-V^2-K}{D},\qquad
 \left(\frac{1-V^2}{D}\right)_u=4V\mathcal F
                                  =-2\mathcal F D_u
 \label{rf:transient-identity}
\end{equation}
follow from \eqref{rf:pair-system}. In the last identity,
$(1-V^2)/D=4M(1-DM)$; differentiating this expression proves the
formula without assigning separate bounds to its factors. We will not
divide its signed integral by an uncertain minimum weight. The
comparisons below use the differential equations directly.

\subsection{The sharp scalar comparison for a constant \texorpdfstring{$K$}{K} floor}

\begin{theorem}[Constant forcing floor with the initial transient retained]
\label{rf:constant-K}
Suppose \eqref{rf:pair-system} holds on a stopped interval with $D>0$,
and suppose $K\ge\kappa$ for a constant $\kappa\le1$. Let $d>0$ and
$v\in\R$ satisfy the two ordered initial conditions
\begin{equation}
 g(0)\ge2\sqrt d,\qquad
 g_u(0)\ge-2v/\sqrt d,\qquad g=2\sqrt D.
 \label{rf:ordered-data}
\end{equation}
Define
\begin{equation}
 d_\kappa(u)=d-2vu+\frac{v^2+\kappa-1}{d}u^2.
 \label{rf:constant-K-polynomial}
\end{equation}
Then $D(u)\ge d_\kappa(u)$ before the first zero of $d_\kappa$ and
before the stipulated pair tube ends. Equality holds for the scalar
comparison with constant $K=\kappa$ and equal initial data. Thus this
comparison retains the full initial transient of that scalar problem.
\end{theorem}

\begin{proof}
Differentiate $g=2\sqrt D$ using \eqref{rf:transient-identity}. The
terms quadratic in $V$ cancel and give
\begin{equation}
 g_{uu}=\frac{16(K-1)}{g^3}
       \ge f_\kappa(g),\qquad
 f_\kappa(g)=\frac{16(\kappa-1)}{g^3}.
 \label{rf:gap-acceleration}
\end{equation}
For $\kappa\le1$, $f_\kappa$ is nondecreasing on $(0,\infty)$.
Direct differentiation of $y=2\sqrt{d_\kappa}$ gives
$y_{uu}=f_\kappa(y)$, $y(0)=2\sqrt d$ and
$y_u(0)=-2v/\sqrt d$, on the initial interval where $d_\kappa>0$.

For completeness, the comparison can be proved on any compact
subinterval on which $g,y>0$. Put $a=g-y$ and $b=g_u-y_u$.
Then $a_u=b$ and
\[
 b_u\ge f_\kappa(g)-f_\kappa(y)=c(u)a,
\]
where the divided difference $c(u)$ is continuous, nonnegative and
bounded, using $f_\kappa'$ at equal arguments. The negative parts obey
$(a_-)_u\le b_-$ and $(b_-)_u\le c(u)a_-$ almost everywhere.
An integrating-factor inequality for $a_-+b_-$, initially zero, makes
both negative parts zero. Hence $g\ge y$, and squaring proves the
assertion. Exhaust the positive comparison interval by such compact
subintervals. No comparison is continued through the first zero of
$d_\kappa$ merely because its quadratic expression has values later.
\end{proof}

The initial conditions in \eqref{rf:ordered-data} are stated for both
position and velocity deliberately. If $D(0)\ge d$ and $V(0)\le v$
with $v\ge0$, they follow at once. A useful alternative uses
$D(0)\ge d$ and $M(0)\ge b\ge0$. Indeed,
\[
 g_u(0)=-\frac2{\sqrt{D(0)}}+4M(0)\sqrt{D(0)}
       \ge-\frac2{\sqrt d}+4b\sqrt d,
\]
because the last expression is increasing in $D$ and $M$ on that
rectangle. Thus one may take $v=1-2db$, without requiring this $v$ to
be nonnegative. These ordering checks also identify the favorable
endpoints when screening an uncertain anchor.

The equality claim is sharpness within the stated scalar differential
comparison, not a claim that every constant-$K$ scalar trajectory is a
canonical divisor. A strictly positive minimum of $d_\kappa-d_{\rm stop}$
on a prescribed finite interval excludes a lower stop, once all other
tube exits have been independently excluded. Conversely, a comparison
that reaches zero before the requested endpoint cannot certify that
endpoint by this theorem. It does not establish an actual collision.

\subsection{A constant physical forcing floor}

A constant floor for $\mathcal F$ is a different input from a constant
floor for $K=2D^2\mathcal F$. In particular, a small $K$ near a small
gap is compatible with a positive physical forcing. The regular
$D/M$ equations retain this distinction without division by $D$.

\begin{theorem}[Regular Riccati and gap comparison]
\label{rf:forcing-comparison}
Suppose \eqref{rf:pair-system} holds on a stopped positive-$D$ interval,
with complete signed $\mathcal F\ge f\ge0$, $D(0)\ge d>0$ and
$M(0)\ge b\ge0$. For $f>0$ put $\lambda=\sqrt{2f}$ and
\begin{equation}
 S(u)=\frac{\sinh(\lambda u)}{\lambda},\qquad
 W(u)=\cosh(\lambda u)+2bS(u),\qquad
 m(u)=\frac{W'(u)}{2W(u)}.
 \label{rf:model-functions}
\end{equation}
For $f=0$ use the continuous limits $S(u)=u$,
$W(u)=1+2bu$ and $m(u)=b/(1+2bu)$. In either case define
\begin{equation}
 d_f(u)=W(u)\bigl[dW(u)-2S(u)\bigr].
 \label{rf:physical-forcing-model}
\end{equation}
Then $M(u)\ge m(u)\ge0$ and $D(u)\ge d_f(u)$ throughout the stopped
interval. The direct regular $D/M$ comparison given the same
$(d,b,f)$ produces exactly these same functions.
\end{theorem}

\begin{proof}
For $u\ge0$, $W\ge1$ and $W'\ge0$. For $f>0$, $W''=2fW$;
the same formula holds at $f=0$ with $W''=0$. Therefore
\[
 m_u=\frac{W''W-(W')^2}{2W^2}=f-2m^2,
 \qquad m(0)=b.
\]
The complete forcing inequality gives
\[
 (M-m)_u+2(M+m)(M-m)\ge0.
\]
On each compact stopped subinterval an integrating factor is positive,
irrespective of the sign of the continuous coefficient. Since
$M(0)-m(0)\ge0$, this proves $M\ge m\ge0$; no exterior-summand
positivity was used.

The Wronskian identity $S'W-SW'=1$ follows directly from the hyperbolic
formulas, and is also valid for $f=0$. Consequently
\[
 (S/W)'=W^{-2},\qquad
 d_f=W^2\left[d-2\int_0^u W(s)^{-2}\dd s\right],\qquad
 (d_f)_u=-2+4m d_f.
\]
Subtracting from the actual equation yields
\[
 (D-d_f)_u-4m(D-d_f)=4D(M-m)\ge0
\]
while $D>0$. A second integrating factor proves $D\ge d_f$.
Both integrations use the regular pair system itself, which proves the
same-information statement as well. No monotonicity of the actual
$D$, and no sign of its actual $V$, is assumed.
\end{proof}

The formula covers the three Riccati regimes without a singular case
distinction. For $f>0$ put $a=\sqrt{f/2}$. If $b>a$, then $m$ decreases
to $a$; if $b=a$, it is constant; if $0\le b<a$, it increases to $a$.
These assertions follow from $m_u=2(a^2-m^2)$ and uniqueness, or from
\eqref{rf:model-functions}. For $f=0$, $m$ decreases to zero when
$b>0$ and is identically zero when $b=0$. In particular,
\begin{equation}
 d_0(u)=(1+2bu)\bigl[d(1+2bu)-2u\bigr],
 \label{rf:zero-forcing-limit}
\end{equation}
recovering the earlier initial-interaction comparison exactly.

\begin{corollary}[A whole-interval lower-stop test]
\label{rf:forcing-stop}
In addition to the hypotheses of Theorem~\ref{rf:forcing-comparison},
suppose an independently cleared count-two contour and unconditional
coordinate enclosures exclude every other exit on $0\le u\le\Delta$.
Acquire $\mathcal F_\Gamma\ge f$ on the corresponding stopped trial
region $D\ge d_{\rm stop}>0$. If
\begin{equation}
 \min_{0\le u\le\Delta}d_f(u)>d_{\rm stop},
 \label{rf:stop-margin}
\end{equation}
then the actual pair stays above that lower stop throughout the interval.
Positivity of $d_f(\Delta)$ alone always implies positivity of $d_f$
on $[0,\Delta]$. The stronger endpoint test
$d_f(\Delta)>d_{\rm stop}$ implies \eqref{rf:stop-margin} whenever
$b\ge\sqrt{f/2}$ and $4bd<2$.
\end{corollary}

\begin{proof}
At a hypothetical first lower exit, the actual coordinates have stayed
in the acquired trial region. Theorem~\ref{rf:forcing-comparison}
therefore gives $D\ge d_f>d_{\rm stop}$ there, a contradiction.
The normalized model satisfies
\[
 \frac{d_f(u)}{W(u)^2}=d-2S(u)/W(u),\qquad
 \left(\frac{d_f}{W^2}\right)'=-2/W^2<0.
\]
Thus a positive endpoint makes the normalized model, and hence $d_f$,
positive at every earlier time. For the stated stronger endpoint test,
$0\le m\le b$. The positive model cannot cross its initial upper
barrier $d$, since at such a crossing its derivative is at most
$-2+4bd<0$. While positive it consequently satisfies
$(d_f)_u\le-2+4bd<0$, so its minimum is its endpoint. In the other
Riccati regimes this last monotonicity is not asserted; use the actual
scalar minimum in \eqref{rf:stop-margin} instead.
\end{proof}

The count and coordinate assumptions in this corollary are the same
independent gates as in Theorem~\ref{cmp:stop}. In particular, coordinate
acquisition may permit negative $D$; the desired positive trajectory is
not an input to it. For a finite family of pair charts the margin must
hold for every covered label and every chart/boundary exit. Passing only
one initially minimizing label does not protect an evolving block
minimum. A regular merging-pair chart also does not regularize its
neighboring edge charts.

\subsection{Complete source acquisition and its error budget}

The trial trace in Proposition~\ref{cmp:trial} gives the combined
physical forcing directly:
\begin{align}
 \mathcal F_\Gamma(t;\widehat c,\widehat D)
 &=-\frac1{2\pi\ii}\oint_\Gamma B(t,z)
       \left(\frac{12}{Q^2}+\frac{16\widehat D}{Q^3}\right)\dd z,
       \notag\\
 &=-\frac1{2\pi\ii}\oint_\Gamma\mathcal L(t,z)Q'(z)
       \left(\frac{24}{Q^3}+\frac{48\widehat D}{Q^4}\right)\dd z,
 \label{rf:source-kernel}
\end{align}
where the second identity requires the matching cleared annular
logarithm of \eqref{cmp:log-kernel}. All interior rational contributions
cancel as in Section~\ref{sec:compact}; every exterior occurrence and
its integer multiplicity remain. Conjugation makes the value real, but
does not make the individual exterior summands nonnegative.

Use the geometry of \eqref{cmp:kernel-errors}: $0\le C<r$,
$|\widehat D|\le D_{\max}$ and $q=(r-C)^2-D_{\max}>0$.
Uniform source errors on the circle give
\begin{align}
 E_{\mathcal F,B}
 &\le re_B\left(\frac{12}{q^2}+\frac{16D_{\max}}{q^3}\right),
       \notag\\
 E_{\mathcal F,\mathcal L}
 &\le2r(r+C)e_{\mathcal L}
       \left(\frac{24}{q^3}+\frac{48D_{\max}}{q^4}\right).
 \label{rf:source-errors}
\end{align}
The quotient estimate \eqref{cmp:quotient-error} pays $e_B$ from
field/first-jet errors and a nonzero contour denominator. These inequalities
follow by taking the modulus of the combined kernel and integrating the
contour length. They need no reciprocal target gap.

If trial coordinates are approximated rather than enclosed as parameters,
suppose their errors are at most $e_c,e_D$ and all connecting parameter
segments retain the same $q$. With $B_{\max}$ bounding the exact field
$|B|$ on the contour, the additional error is at most
\begin{equation}
 \begin{split}
 E_{\rm par}\le rB_{\max}\bigg[&
 \left(\frac{48(r+C)}{q^3}
             +\frac{96D_{\max}(r+C)}{q^4}\right)e_c\\
 &+\left(\frac{40}{q^3}+\frac{48D_{\max}}{q^4}\right)e_D\bigg].
 \end{split}
 \label{rf:parameter-error}
\end{equation}
To verify the constants, differentiate
$J=12Q^{-2}+16\widehat DQ^{-3}$, using
$Q_{\widehat c}=-2(z-\widehat c)$ and $Q_{\widehat D}=-1$:
\[
 J_{\widehat c}=48(z-\widehat c)Q^{-3}
                    +96\widehat D(z-\widehat c)Q^{-4},\qquad
 J_{\widehat D}=40Q^{-3}+48\widehat DQ^{-4}.
\]
Integrate the parameter derivatives along the joining segments and then
around the contour. Evaluating all parameters by a rigorous enclosure
can instead include this payment directly, without adding it twice.

The remaining requirements are precisely those of
\eqref{cmp:trapezoid} and its following source ledger: quadrature on a
cleared angular strip, complete theta/xi and integration tails, arithmetic,
reconstruction, spatial uncertainty and a genuine whole-time remainder.
Heat identities for time derivatives do not remove the remainder.
A valid lower acquisition has the form
\[
 \inf\widehat{\mathcal F}_\Gamma
 -E_{\rm source}-E_{\rm par}-E_{\rm quad}
 -E_{\rm arithmetic}-E_{\rm time}-E_{\rm reconstruction}\ge f,
\]
with a nonoverlapping accounting of errors. It must cover the actual
stopped coordinates and all needed time cells. A bound for
$|\mathcal F|$ is not this signed lower estimate. A contour representation
does not leave an unpaid zero-tail cutoff once the complete source error
has been paid; a finite source implementation still has to pay that error.

Replacing the contour by root jets requires a lower bound for $|H_z|$
and root-location/reciprocal errors. Converting a $D$ error to a gap error
on $D\ge d_{\rm stop}$ costs at most $1/\sqrt{d_{\rm stop}}$ times
that error. Such denominators are not implicit in
\eqref{rf:source-errors}. Negative-time source bounds likewise need
their own domain verification as in \eqref{cmp:negative-gaussian}.
All these inputs must be given equally to direct regular $D/M$ evolution
or direct contour-discriminant acquisition. Theorem~\ref{rf:forcing-comparison}
claims no additional margin or cost advantage over those methods.

\subsection{Forward reanchoring and the initial-interaction horizon}

Return to forward physical time $t$. This subsection concerns one
specific consumer: at a starting slice, discard the positive forcing
after retaining the exact initial $D,M$, and use
\eqref{rf:zero-forcing-limit} backward. On $D>0$, $M\ge0$ and $DM<1$
its strict positive-gap horizon is
\begin{equation}
 h(t)=\frac{D(t)}{2(1-D(t)M(t))}.
 \label{rf:horizon-definition}
\end{equation}
This is an allowance, not the actual time remaining until a collision.
The consumer remains conditional on its backward source/sign regime and
coverage; an algebraic horizon alone does not acquire that regime.

\begin{theorem}[Exact evolution of the initial-$M$ horizon]
\label{rf:horizon}
On a forward local interval suppose the complete pair equations hold,
$D>0$, $M\ge0$ and $\mathcal F\ge0$. Assume $1-D(0)M(0)>0$.
Then $1-DM>0$ throughout that interval, and
\begin{equation}
 h_t=1-2\mathcal Fh^2,\qquad
 h(s)-s=h(0)-2\int_0^s\mathcal F(t)h(t)^2\dd t\le h(0).
 \label{rf:horizon-identity}
\end{equation}
Consequently, forwarding to $s$ and applying the unchanged
complete-initial-$M$ consumer there cannot reach an earlier endpoint
than the same exact-input consumer at time zero.
\end{theorem}

\begin{proof}
The forward equations are $D_t=2-4DM$ and
$M_t=2M^2-\mathcal F$. Put $y=1-DM$. Direct differentiation gives
\[
 y_t=-D_tM-DM_t=-2My+D\mathcal F.
\]
On every compact forward subinterval, multiplying by
$\exp(2\int_0^tM)$ proves $y(t)>0$ from $y(0)>0$, because
$D\mathcal F\ge0$. Thus the denominator of $h=D/(2y)$ does not
cross zero. Differentiating the quotient now gives
\[
 h_t=\frac{D_ty-Dy_t}{2y^2}
     =\frac{2y^2-D^2\mathcal F}{2y^2}
     =1-2\mathcal Fh^2.
\]
Integration proves \eqref{rf:horizon-identity}. The strict backward
allowance from time $s$ reaches only times greater than $s-h(s)$,
and the displayed identity gives $s-h(s)\ge-h(0)$.
\end{proof}

The hypothesis $M\ge0$ here is explicit throughout the forward
interval. It is not inferred from $\mathcal F\ge0$: at $M=0$ the
equation $M_t=-\mathcal F$ can point toward negative values. A complete
all-real adjacent-pair regime supplies both signs, while a local tube
with possibly nonreal exterior needs them separately. The raw quotient
identity holds more generally whenever $D>0$ and $DM\ne1$; the theorem
keeps the domain of the inherited nonnegative-initial-$M$ consumer.

For any persistent finite family of eligible labels, apply
\eqref{rf:horizon-identity} to each label. Every $h_j(t)-t$ is
nonincreasing, so is their minimum on each fixed block, and the sets of
block indices whose allowance exceeds a fixed deadline threshold can
only shrink. Repeated exact reanchoring at increasing forward times
within this same domain does not change that conclusion. A changed
minimizing label creates no improvement, since the pointwise inequalities
hold for all labels simultaneously. Coverage and outer-boundary exits
are still separate hypotheses.

The same ceiling applies to a nonnegative two-neighbor lower estimate
$b_3(s)\le M(s)$ only when that inequality is actually paid. For fixed
$D$, the horizon $D/[2(1-Db)]$ is increasing in $b$ on $Db<1$, so the
corresponding horizon is at most $h(s)$. In a complete all-real regime,
positive omission supplies this lower estimate. If the remaining exterior
may be nonreal, nonnegative complete $M$ alone does not make a positive
real-neighbor subsum a lower bound. Its omitted exterior must be paid.
This comparison is with the \emph{complete exact} initial-$M$ allowance
at time zero; it does not rule out improving a weaker or uncertain
initial estimate.

The quadratic heat collision has $\mathcal F=0$ and saturates the
round-trip identity. For stationary cosine the divisor does not move,
$h$ is constant and $h-t$ decreases: loss of an allowance is therefore
not loss of the actual gap. The exclusion is restricted to the unchanged
initial-$M$ consumer under the stated sign and common-label assumptions.
Signed forcing, the retained-$\mathcal F$ comparison above, time-dependent
joint evolution, and additional source information or acquisition-cost
gains are outside it. The generic identities apply to appropriate
changed-source controls; canonical use still requires the complete
unchanged source and every count, time and error premise.
